\section{The AreaSection Algorithm}
\label{sec:algorithm}

\subsection{Setup and Notation}

Let $G = (V, E, w)$ be the state's census-tract adjacency graph with vertex
weights $\mathbf{p} \in \mathbb{Z}_{>0}^{|V|}$ (population) and edge weights
$w(e) > 0$ (shared boundary length in metres, from TIGER).
Let $k$ be the number of congressional districts.
Let $\mathbf{a} \in \mathbb{R}_{>0}^{|V|}$ be per-vertex land area (ALAND
in square metres from TIGER, scaled to hectares).

\begin{definition}[Dual-constrained bisection]
A \emph{dual-constrained bisection} of $G$ into parts $(S, V \setminus S)$
with parameters $(p_L, \alpha)$ satisfies:
\begin{align}
  \frac{\sum_{v \in S} p_v}{\sum_{v \in V} p_v} &\in \left[\frac{p_L}{1.001},\; 1.001 \cdot p_L\right]
    & \text{(population balance)} \\
  \frac{\sum_{v \in S} a_v}{\sum_{v \in V} a_v} &\in \left[\frac{\alpha}{1.10},\; 1.10 \cdot \alpha\right]
    & \text{(area balance)}
\end{align}
where $p_L$ is the target population fraction for $S$ and $\alpha = 0.5$
(50\,\% land area target).
\end{definition}

The population tolerance is tight (0.1\,\%, corresponding to METIS
$\mathtt{ubvec}[0] = 1.001$).
The area tolerance is loose (10\,\%, $\mathtt{ubvec}[1] = 1.10$) to allow
METIS to find a good population-balanced partition while making a best-effort
attempt at area balance.
When the constraints conflict --- which occurs in states with highly
concentrated populations --- METIS prioritises the tighter population
constraint.

\subsection{Lorenz Feasibility Pre-filter}

Before trying any METIS calls, AreaSection computes the
\emph{population-area Lorenz curve} to identify infeasible ratios.

\begin{definition}[Population-area Lorenz curve]
Sort tracts by density $\rho_v = p_v / a_v$ in descending order.
The population-area Lorenz curve $L : [0,1] \to [0,1]$ maps cumulative area
fraction $x$ to cumulative population fraction $L(x)$:
\[
  L(x) = \frac{\sum_{v : \mathrm{rank}(v) \leq x|V|} p_v}{\sum_v p_v}
\]
where tracts are ordered from densest to least dense.
\end{definition}

The Lorenz curve provides lower and upper bounds on the achievable area fraction
for a given population fraction:

\begin{proposition}[Feasibility bounds]
For a target population fraction $p$, the set of achievable non-contiguous
area fractions is the interval
\[
  [L^{-1}(p),\; 1 - L^{-1}(1-p)]
\]
where $L^{-1}$ is the inverse Lorenz function (minimum area fraction to
achieve population fraction $p$).
A ratio $i:(k-i)$ with $p = i/k$ is \emph{Lorenz-infeasible} if this interval
does not overlap $[0.5/1.10,\; 0.5 \times 1.10] = [0.455,\; 0.55]$.
\end{proposition}

\begin{proof}
$L^{-1}(p)$ is the minimum area fraction achievable for population fraction $p$
(greedy dense-first selection). $1 - L^{-1}(1-p)$ is the maximum area fraction
achievable (greedy sparse-first, equivalently leaving the densest $1-p$
fraction to the other side). Contiguity tightens these bounds further, so
Lorenz-infeasibility is a necessary condition for infeasibility under the
contiguity constraint.
\end{proof}

\begin{remark}[Necessary, not sufficient]
Proposition~\ref{prop:lorenz-feasibility} establishes a \emph{necessary} but
not \emph{sufficient} condition for feasibility.
A ratio may be \emph{Lorenz-feasible} (the interval overlaps the allowed area
window) yet infeasible for the \emph{contiguous} partitioning problem that METIS
solves: geographic chokepoints, elongated state shapes, and connected-subgraph
requirements all tighten the achievable set beyond the non-contiguous lower bound.
We use ``Lorenz-feasible ratio'' throughout to mean a ratio passing this filter,
not a ratio for which METIS is guaranteed to succeed.
\end{remark}

\noindent\textbf{Natural ratio.}
The Lorenz curve also identifies the \emph{natural population fraction}: the
value $p^* = L(0.5)$, i.e., the population fraction contained in the densest
50\,\% of land area.
For a uniform-density state, $p^* = 0.5$ (equal split is natural).
For Wisconsin, $p^* = 0.932$: the densest half of Wisconsin's land holds 93.2\,\%
of the population, confirming that the 50/50 area constraint will force a
near-equal population split.

\subsection{Ratio Scanning and Normalisation}

At the first bisection level (the ``natural ratio'' level), AreaSection tries
all feasible ratios $i:(k-i)$ for $i = 1, \ldots, \lfloor k/2 \rfloor$,
subject to Lorenz feasibility.
For each ratio, it runs $N$ METIS seeds and records the minimum edge-cut.
The winning ratio minimises the \emph{isoperimetrically normalised edge-cut}:
\[
  \mathrm{EC}_{\text{norm}}(i:k-i) = \frac{\mathrm{EC}(i:k-i)}{\sqrt{\min(i,\, k-i)}}
\]

The normalisation is identical to GeoSection (B.8).
Its purpose is to prevent the caterpillar pathology: without normalisation,
1:k-1 always wins because a 1-district boundary is cheaper than a k/2:k/2
boundary regardless of compactness.
The $\sqrt{\min(i,k-i)}$ factor is the isoperimetric correction for a convex
region, ensuring that both large and small splits compete on equal footing.

\subsection{METIS Dual-Constraint Implementation}

Each METIS call constructs a graph with $\mathtt{ncon}=2$ vertex weights
(population, area in hectares) and a 4-element $\mathtt{tpwgts}$ vector:
\[
  \mathtt{tpwgts} = \left[\frac{i}{k},\; 0.5,\; \frac{k-i}{k},\; 0.5\right]
\]
specifying that partition 0 targets $i/k$ of the population and 50\,\% of
the land area, and partition 1 targets the remainder.
The $\mathtt{ubvec} = [1.001, 1.10]$ specifies \emph{soft} imbalance tolerances:
METIS \emph{targets} 0.1\,\% population balance and 10\,\% area balance but may
exceed the area tolerance when the two constraints conflict.
Because population balance uses a tighter tolerance ($\mathtt{ubvec}[0] = 1.001$),
METIS prioritises population over area when they cannot be simultaneously
satisfied.
The correct interpretation is that AreaSection \emph{penalises} area imbalance
and \emph{targets} geographic balance; it does not \emph{enforce} a hard area
constraint.

Area values are scaled from square metres to hectares ($\div 10{,}000$) to
remain within METIS's 32-bit integer range.
Wisconsin's largest census tract has 138{,}992 hectares, well within the
$2^{31} - 1 \approx 2.1 \times 10^9$ limit.

\paragraph{Contiguity.}
The bisection path does not set METIS's \texttt{-contig} option at the
$\mathtt{ncon}=1$ recursive levels.
For planar census tract graphs, METIS empirically produces contiguous partitions
without the flag, as established by the metis\_format.rs documentation in the
reference implementation.
Contiguity is not \emph{algorithmically} guaranteed; it is \emph{empirically}
reliable for US state adjacency graphs.
States where a non-contiguous partition is detected (rare) trigger a post-hoc
repair that reassigns isolated components to the adjacent partition.

\subsection{Recursive Structure}

After the top-level natural ratio is determined, AreaSection recurses into
each half using GeoSection's standard ratio-scanning bisection with
\emph{population-only} vertex weights ($\mathtt{ncon}=1$).
The area constraint applies only at the first bisection level.
This design reflects the legal theory: the natural ratio is the geographic
structure of the state, and subsequent levels are standard minimum-edge-cut
bisection within each geographic half.
The geographic balance property is therefore a first-level property of the
bisection tree, not a guarantee on individual district areas.

The full algorithm is:
\begin{enumerate}
  \item Compute Lorenz curve; determine natural ratio $p^*$ and feasible ratios.
  \item For each feasible ratio $i:(k-i)$, run $N$ dual-constraint METIS seeds;
    record minimum $\mathrm{EC}_{\text{norm}}$.
  \item Select ratio with minimum $\mathrm{EC}_{\text{norm}}$ as the natural ratio.
  \item Recurse on each half with GeoSection (ncon=1, ratio scanning).
\end{enumerate}

\subsection{Computational Cost}

AreaSection adds one dual-constraint METIS call per ratio per seed at the
first level.
For a state with $k$ districts and $N$ seeds per ratio, this is
$\lfloor k/2 \rfloor \times N$ additional METIS calls compared to standard
bisection.
All subsequent recursion levels use the identical ncon=1 GeoSection path.

In practice, Lorenz pre-filtering eliminates 0--2 ratios per state, providing
modest speedup.
For Wisconsin (8 districts, 50 seeds), the first-level scan runs
$4 \times 50 = 200$ dual-constraint calls; subsequent levels run the same
200 ncon=1 calls as GeoSection.
Wall-clock time is approximately $2\times$ GeoSection for the first level only.
